\documentclass[12pt]{article}
\usepackage{amsmath}
\usepackage{amssymb}
\usepackage{geometry}
\usepackage{hyperref}
\geometry{a4paper, margin=1in}

\title{Linear Independence: Comprehensive Problem Sets \& Solutions}
\author{Based on Homework 1.7 Framework}
\date{}

\begin{document}

\maketitle

\tableofcontents
\newpage

\section{Concept Checklist \& Question Mapping}
This checklist serves as a comprehensive inventory of all linear independence concepts, tricks, and problem types extracted from the source material. It maps each concept to the relevant questions in the problem set to ensure total coverage.

\begin{itemize}
    \item \textbf{C1: Algorithmic Reduction (Row Reduction / Gaussian Elimination)} 
    \newline \textit{Concept:} Setting up an augmented/coefficient matrix and reducing it to determine if every column has a pivot (no free variables).
    \newline \textit{Questions:} 1, 2, 3
    
    \item \textbf{C2: Independence of Matrix Columns}
    \newline \textit{Concept:} Analyzing $A\mathbf{x} = \mathbf{0}$ for specific dimensions of matrices to check for pivot columns.
    \newline \textit{Questions:} 4, 5
    
    \item \textbf{C3: Unknown Variables and Spans}
    \newline \textit{Concept:} Carrying an unknown variable $h$ through row operations or determinants to find boundaries of linear dependence/independence.
    \newline \textit{Questions:} 6, 7, 8, 9
    
    \item \textbf{C4: Visual Inspection - Oversize Sets (Theorem 8)}
    \newline \textit{Concept:} Identifying that $p > n$ (more vectors than dimensions) strictly forces a linearly dependent set.
    \newline \textit{Questions:} 10, 11
    
    \item \textbf{C5: Visual Inspection - Zero Vector Rule (Theorem 9)}
    \newline \textit{Concept:} Identifying that any set containing the zero vector $\mathbf{0}$ is automatically linearly dependent.
    \newline \textit{Questions:} 12, 13
    
    \item \textbf{C6: Visual Inspection - Two-Vector Rule}
    \newline \textit{Concept:} For exactly two vectors, determining independence by checking if one is a scalar multiple of the other.
    \newline \textit{Questions:} 14, 15
    
    \item \textbf{C7: Conceptual True/False}
    \newline \textit{Concept:} Testing theoretical understanding of the trivial solution, spans, and matrix dimensions.
    \newline \textit{Questions:} 16, 17, 18
    
    \item \textbf{C8: Echelon Form Analysis}
    \newline \textit{Concept:} Working backwards from the definition of independence to determine the required geometric shape and pivot count of matrices.
    \newline \textit{Questions:} 19, 20
    
    \item \textbf{C9: Advanced Applications - Linear Independence of Functions (Wronskian)}
    \newline \textit{Concept:} Using the Wronskian determinant to test functions (which form a vector space).
    \newline \textit{Questions:} 21, 22, 23
    
    \item \textbf{C10: Advanced Applications - Independence via Determinants}
    \newline \textit{Concept:} For square $n \times n$ matrices, using $\det(A) \neq 0$ as a fast-track test for independence.
    \newline \textit{Questions:} 24, 25
    
    \item \textbf{C11: Real-World Modeling \& Quantitative Finance}
    \newline \textit{Concept:} Setting up replicating portfolios and identifying arbitrage opportunities via redundant (dependent) payoff vectors.
    \newline \textit{Questions:} 26, 27
    
    \item \textbf{C12: Diagnostic Testing ("Find the Flaw")}
    \newline \textit{Concept:} Identifying subtle conceptual or algebraic errors in flawed arguments regarding independence.
    \newline \textit{Questions:} 28, 29, 30
\end{itemize}

\newpage

\section{Problem Sets}

\subsection*{Type 1: Algorithmic Reduction}
\textbf{Problem 1:} Determine if the vectors are linearly independent. Justify your answer using row reduction.
$$ \mathbf{v}_1 = \begin{bmatrix} 1 \\ 2 \\ 1 \end{bmatrix}, \quad \mathbf{v}_2 = \begin{bmatrix} 2 \\ 5 \\ 0 \end{bmatrix}, \quad \mathbf{v}_3 = \begin{bmatrix} 3 \\ 3 \\ 8 \end{bmatrix} $$

\textbf{Problem 2:} Determine if the vectors are linearly independent.
$$ \mathbf{v}_1 = \begin{bmatrix} 1 \\ 1 \\ 1 \end{bmatrix}, \quad \mathbf{v}_2 = \begin{bmatrix} 2 \\ 3 \\ 4 \end{bmatrix}, \quad \mathbf{v}_3 = \begin{bmatrix} 4 \\ 5 \\ 6 \end{bmatrix} $$

\textbf{Problem 3:} Determine if the set of vectors in $\mathbb{R}^4$ is linearly independent.
$$ \mathbf{v}_1 = \begin{bmatrix} 1 \\ 0 \\ 2 \\ 0 \end{bmatrix}, \quad \mathbf{v}_2 = \begin{bmatrix} 0 \\ 1 \\ -1 \\ 1 \end{bmatrix}, \quad \mathbf{v}_3 = \begin{bmatrix} 1 \\ 1 \\ 1 \\ 2 \end{bmatrix} $$

\subsection*{Type 2: Independence of Matrix Columns}
\textbf{Problem 4:} Determine if the columns of matrix $A$ form a linearly independent set.
$$ A = \begin{bmatrix} 2 & -1 & 4 \\ -4 & 2 & -8 \\ 1 & 1 & 1 \end{bmatrix} $$

\textbf{Problem 5:} Evaluate whether the homogeneous equation $A\mathbf{x} = \mathbf{0}$ has only the trivial solution, given:
$$ A = \begin{bmatrix} 1 & 2 & 3 & 4 \\ 0 & 5 & 6 & 7 \\ 0 & 0 & 8 & 9 \\ 0 & 0 & 0 & 2 \end{bmatrix} $$

\subsection*{Type 3: Unknown Variables and Spans}
\textbf{Problem 6:} Find the value of $h$ for which the vectors are linearly dependent.
$$ \mathbf{v}_1 = \begin{bmatrix} 1 \\ 2 \\ 3 \end{bmatrix}, \quad \mathbf{v}_2 = \begin{bmatrix} -1 \\ 0 \\ 2 \end{bmatrix}, \quad \mathbf{v}_3 = \begin{bmatrix} 1 \\ 6 \\ h \end{bmatrix} $$

\textbf{Problem 7:} Find the value(s) of $h$ for which the columns of the matrix are linearly \textit{independent}.
$$ \begin{bmatrix} 2 & 0 & 6 \\ -4 & 3 & 0 \\ 1 & -1 & h \end{bmatrix} $$

\textbf{Problem 8:} Find the values of $h$ for which the vectors $\mathbf{v}_1 = \begin{bmatrix} h \\ 4 \end{bmatrix}$ and $\mathbf{v}_2 = \begin{bmatrix} 9 \\ h \end{bmatrix}$ are linearly independent. 

\textbf{Problem 9:} Let $\mathbf{v}_1 = \begin{bmatrix} 1 \\ -2 \\ 1 \end{bmatrix}$, $\mathbf{v}_2 = \begin{bmatrix} -2 \\ 4 \\ -2 \end{bmatrix}$, and $\mathbf{v}_3 = \begin{bmatrix} 3 \\ 1 \\ h \end{bmatrix}$. 
\newline (a) For what values of $h$ is $\mathbf{v}_3$ in Span$\{\mathbf{v}_1, \mathbf{v}_2\}$?
\newline (b) For what values of $h$ is the set $\{\mathbf{v}_1, \mathbf{v}_2, \mathbf{v}_3\}$ linearly dependent?

\subsection*{Type 4 \& 5: Visual Inspection (Theorems 8, 9, \& Two-Vector Rule)}
\textbf{Problem 10:} Determine by inspection if the vectors are independent: 5 vectors in $\mathbb{R}^4$. Justify using a theorem.

\textbf{Problem 11:} Determine by inspection if the vectors are independent:
$$ \begin{bmatrix} 1 \\ 2 \end{bmatrix}, \begin{bmatrix} 0 \\ -1 \end{bmatrix}, \begin{bmatrix} 3 \\ 3 \end{bmatrix}, \begin{bmatrix} 4 \\ -5 \end{bmatrix} $$

\textbf{Problem 12:} Determine by inspection if the set is linearly independent:
$$ \mathbf{v}_1 = \begin{bmatrix} 1 \\ 2 \\ 9 \end{bmatrix}, \quad \mathbf{v}_2 = \begin{bmatrix} 0 \\ 0 \\ 0 \end{bmatrix}, \quad \mathbf{v}_3 = \begin{bmatrix} 3 \\ 4 \\ -2 \end{bmatrix} $$

\textbf{Problem 13:} Without row reduction, determine if the columns of the matrix are independent:
$$ \begin{bmatrix} 3 & -1 & 0 & 4 \\ 2 & 5 & 0 & 1 \\ -7 & 2 & 0 & 3 \end{bmatrix} $$

\textbf{Problem 14:} Determine by inspection if $\mathbf{v}_1 = \begin{bmatrix} 3 \\ -6 \\ 9 \end{bmatrix}$ and $\mathbf{v}_2 = \begin{bmatrix} -1 \\ 2 \\ -3 \end{bmatrix}$ are independent.

\textbf{Problem 15:} Determine by inspection if $\mathbf{v}_1 = \begin{bmatrix} 2 \\ 5 \end{bmatrix}$ and $\mathbf{v}_2 = \begin{bmatrix} 5 \\ 2 \end{bmatrix}$ are independent.

\subsection*{Type 6: Conceptual True/False}
\textbf{Problem 16:} True or False: If the equation $A\mathbf{x} = \mathbf{0}$ has the trivial solution, then the columns of $A$ are linearly independent.

\textbf{Problem 17:} True or False: If a set contains fewer vectors than there are entries in each vector ($p < n$), the set must be linearly independent.

\textbf{Problem 18:} True or False: A set of 3 vectors in $\mathbb{R}^5$ is linearly independent if the corresponding $5 \times 3$ matrix has 3 pivot columns.

\subsection*{Type 7: Echelon Form Analysis}
\textbf{Problem 19:} Describe the possible echelon form(s) of a $4 \times 2$ matrix with linearly independent columns. Use $\blacksquare$ for a non-zero pivot and $*$ for any real number.

\textbf{Problem 20:} How many pivot columns must a $6 \times 4$ matrix have if its columns are linearly independent? Why?

\subsection*{Type 8: Advanced - The Wronskian and Determinants}
\textbf{Problem 21:} Use the Wronskian to determine if the functions $f(x) = e^{2x}$ and $g(x) = e^{3x}$ are linearly independent.

\textbf{Problem 22:} Use the Wronskian to determine if $f(x) = 1$, $g(x) = x$, and $h(x) = e^x$ are linearly independent.

\textbf{Problem 23:} Are the functions $f(x) = \cos^2(x)$, $g(x) = \sin^2(x)$, and $h(x) = 1$ linearly independent? (Hint: Consider algebraic relationships before calculating).

\textbf{Problem 24:} Use the determinant method to test if the columns of $A$ are independent:
$$ A = \begin{bmatrix} 1 & 2 & 3 \\ 0 & 1 & 4 \\ 0 & 0 & 5 \end{bmatrix} $$

\textbf{Problem 25:} For what values of $k$ are the vectors linearly independent?
$$ \mathbf{v}_1 = \begin{bmatrix} k \\ 1 \\ 0 \end{bmatrix}, \quad \mathbf{v}_2 = \begin{bmatrix} 0 \\ k \\ 1 \end{bmatrix}, \quad \mathbf{v}_3 = \begin{bmatrix} 1 \\ 0 \\ k \end{bmatrix} $$

\subsection*{Type 9: Real-World Modeling (Quantitative Finance)}
\textbf{Problem 26:} A market has 3 economic states. The payoff vectors for a Stock (S), Bond (B), and a Call Option (C) are $S = \begin{bmatrix} 10 \\ 20 \\ 30 \end{bmatrix}$, $B = \begin{bmatrix} 10 \\ 10 \\ 10 \end{bmatrix}$, and $C = \begin{bmatrix} 0 \\ 10 \\ 20 \end{bmatrix}$. Set up the matrix equation to check if the Option can be perfectly replicated by the Stock and Bond, and determine if the set $\{S, B, C\}$ is independent.

\textbf{Problem 27:} An analyst evaluates 3 assets in a simplified market with only 2 future states. Must there be a redundant asset? Justify your answer using linear algebra theorems.

\subsection*{Type 10: Diagnostic - "Find the Flaw"}
\textbf{Problem 28:} \textit{Find the Flaw:} A student is given 3 vectors in $\mathbb{R}^5$. They state: "Because $5 > 3$, there are more dimensions than vectors. Therefore, by Theorem 8, the set is automatically linearly dependent." Identify and correct the flaw.

\textbf{Problem 29:} \textit{Find the Flaw:} A student row reduces a $3 \times 3$ matrix and obtains:
$\begin{bmatrix} 1 & 0 & 2 \\ 0 & 1 & 3 \\ 0 & 0 & 0 \end{bmatrix}$. They conclude: "Since every row that isn't all zeros contains a pivot, the vectors are linearly independent." Identify and correct the error.

\textbf{Problem 30:} \textit{Find the Flaw:} To test if $\mathbf{v}_1, \mathbf{v}_2, \mathbf{v}_3$ are independent, a student evaluates $c_1\mathbf{v}_1 + c_2\mathbf{v}_2 + c_3\mathbf{v}_3 = \mathbf{0}$. They set $c_1=0, c_2=0, c_3=0$, which satisfies the equation. The student concludes: "Since the equation is satisfied when the weights are zero, the vectors are linearly independent." Identify and correct the logical error.

\newpage
\section{Detailed Solutions}

\textbf{Solution 1:} We set up the matrix with the vectors as columns and reduce to echelon form.
$$ \begin{bmatrix} 1 & 2 & 3 \\ 2 & 5 & 3 \\ 1 & 0 & 8 \end{bmatrix} \xrightarrow{\substack{R_2 - 2R_1 \\ R_3 - R_1}} \begin{bmatrix} 1 & 2 & 3 \\ 0 & 1 & -3 \\ 0 & -2 & 5 \end{bmatrix} \xrightarrow{R_3 + 2R_2} \begin{bmatrix} 1 & 2 & 3 \\ 0 & 1 & -3 \\ 0 & 0 & -1 \end{bmatrix} $$
There is a pivot in every column (columns 1, 2, and 3). There are no free variables. \textbf{Final Answer: Linearly independent.}

\textbf{Solution 2:} Set up the matrix and row reduce:
$$ \begin{bmatrix} 1 & 2 & 4 \\ 1 & 3 & 5 \\ 1 & 4 & 6 \end{bmatrix} \xrightarrow{\substack{R_2 - R_1 \\ R_3 - R_1}} \begin{bmatrix} 1 & 2 & 4 \\ 0 & 1 & 1 \\ 0 & 2 & 2 \end{bmatrix} \xrightarrow{R_3 - 2R_2} \begin{bmatrix} 1 & 2 & 4 \\ 0 & 1 & 1 \\ 0 & 0 & 0 \end{bmatrix} $$
Column 3 lacks a pivot position, meaning $x_3$ is a free variable. \textbf{Final Answer: Linearly dependent.}

\textbf{Solution 3:} Set up the $4 \times 3$ matrix and row reduce:
$$ \begin{bmatrix} 1 & 0 & 1 \\ 0 & 1 & 1 \\ 2 & -1 & 1 \\ 0 & 1 & 2 \end{bmatrix} \xrightarrow{R_3 - 2R_1} \begin{bmatrix} 1 & 0 & 1 \\ 0 & 1 & 1 \\ 0 & -1 & -1 \\ 0 & 1 & 2 \end{bmatrix} \xrightarrow{\substack{R_3 + R_2 \\ R_4 - R_2}} \begin{bmatrix} 1 & 0 & 1 \\ 0 & 1 & 1 \\ 0 & 0 & 0 \\ 0 & 0 & 1 \end{bmatrix} \xrightarrow{\text{Swap } R_3, R_4} \begin{bmatrix} 1 & 0 & 1 \\ 0 & 1 & 1 \\ 0 & 0 & 1 \\ 0 & 0 & 0 \end{bmatrix} $$
All 3 columns contain a pivot. No free variables. \textbf{Final Answer: Linearly independent.}

\textbf{Solution 4:} Row reduce matrix $A$:
$$ \begin{bmatrix} 2 & -1 & 4 \\ -4 & 2 & -8 \\ 1 & 1 & 1 \end{bmatrix} \xrightarrow{\text{Swap } R_1, R_3} \begin{bmatrix} 1 & 1 & 1 \\ -4 & 2 & -8 \\ 2 & -1 & 4 \end{bmatrix} \xrightarrow{\substack{R_2 + 4R_1 \\ R_3 - 2R_1}} \begin{bmatrix} 1 & 1 & 1 \\ 0 & 6 & -4 \\ 0 & -3 & 2 \end{bmatrix} \xrightarrow{R_3 + \frac{1}{2}R_2} \begin{bmatrix} 1 & 1 & 1 \\ 0 & 6 & -4 \\ 0 & 0 & 0 \end{bmatrix} $$
Column 3 lacks a pivot. \textbf{Final Answer: The columns are linearly dependent.}

\textbf{Solution 5:} The matrix is already in upper triangular (echelon) form. The main diagonal contains the entries 1, 5, 8, 2. Since none of these are zero, there is a pivot in all 4 columns. \textbf{Final Answer: The columns are linearly independent} (the system only has the trivial solution).

\textbf{Solution 6:} Set up the matrix and isolate $h$ via row operations:
$$ \begin{bmatrix} 1 & -1 & 1 \\ 2 & 0 & 6 \\ 3 & 2 & h \end{bmatrix} \xrightarrow{\substack{R_2 - 2R_1 \\ R_3 - 3R1}} \begin{bmatrix} 1 & -1 & 1 \\ 0 & 2 & 4 \\ 0 & 5 & h-3 \end{bmatrix} \xrightarrow{R_3 - \frac{5}{2}R_2} \begin{bmatrix} 1 & -1 & 1 \\ 0 & 2 & 4 \\ 0 & 0 & (h-3) - 10 \end{bmatrix} = \begin{bmatrix} 1 & -1 & 1 \\ 0 & 2 & 4 \\ 0 & 0 & h-13 \end{bmatrix} $$
For the vectors to be dependent, there must be a free variable, requiring the pivot in column 3 to be zero: $h - 13 = 0 \implies h = 13$. \textbf{Final Answer: Dependent for $h = 13$.}

\textbf{Solution 7:} Set up the matrix:
$$ \begin{bmatrix} 2 & 0 & 6 \\ -4 & 3 & 0 \\ 1 & -1 & h \end{bmatrix} \xrightarrow{\text{Swap } R_1, R_3} \begin{bmatrix} 1 & -1 & h \\ -4 & 3 & 0 \\ 2 & 0 & 6 \end{bmatrix} \xrightarrow{\substack{R_2 + 4R_1 \\ R_3 - 2R_1}} \begin{bmatrix} 1 & -1 & h \\ 0 & -1 & 4h \\ 0 & 2 & 6-2h \end{bmatrix} $$
$$ \xrightarrow{R_3 + 2R_2} \begin{bmatrix} 1 & -1 & h \\ 0 & -1 & 4h \\ 0 & 0 & (6-2h) + 8h \end{bmatrix} = \begin{bmatrix} 1 & -1 & h \\ 0 & -1 & 4h \\ 0 & 0 & 6h+6 \end{bmatrix} $$
For independence, there must be no free variables. The pivot must NOT be zero: $6h + 6 \neq 0 \implies h \neq -1$. \textbf{Final Answer: Independent for all $h \neq -1$.}

\textbf{Solution 8:} Using the determinant trick for a $2 \times 2$ matrix: $A = \begin{bmatrix} h & 9 \\ 4 & h \end{bmatrix}$. 
$\det(A) = h^2 - 36$. For independence, the determinant must be non-zero. $h^2 - 36 \neq 0 \implies h \neq \pm 6$. \textbf{Final Answer: Independent for $h \neq 6$ and $h \neq -6$.}

\textbf{Solution 9:} (a) Notice $\mathbf{v}_2 = -2\mathbf{v}_1$. The span of $\mathbf{v}_1$ and $\mathbf{v}_2$ is just the line defined by $\mathbf{v}_1$. If $\mathbf{v}_3 \in \text{Span}\{\mathbf{v}_1, \mathbf{v}_2\}$, then $\mathbf{v}_3 = c\mathbf{v}_1$. 
$\begin{bmatrix} 3 \\ 1 \\ h \end{bmatrix} = c\begin{bmatrix} 1 \\ -2 \\ 1 \end{bmatrix}$. From row 1, $c = 3$. But from row 2, $1 \neq 3(-2)$. Contradiction. \textbf{Answer (a): No values of $h$.}
(b) Since $\mathbf{v}_2 = -2\mathbf{v}_1$, the subset $\{\mathbf{v}_1, \mathbf{v}_2\}$ is already linearly dependent. The addition of $\mathbf{v}_3$ does not change this. \textbf{Answer (b): The set is dependent for all values of $h$.}

\textbf{Solution 10:} There are $p = 5$ vectors in $n = 4$ dimensions. According to Theorem 8, any set containing more vectors than entries in each vector is linearly dependent. \textbf{Final Answer: Linearly dependent (Theorem 8).}

\textbf{Solution 11:} There are $p = 4$ vectors in $n = 2$ dimensions. $4 > 2$. \textbf{Final Answer: Linearly dependent (Theorem 8).}

\textbf{Solution 12:} The set explicitly contains the zero vector $\mathbf{v}_2 = \mathbf{0}$. According to Theorem 9, any set containing the zero vector is automatically linearly dependent. \textbf{Final Answer: Linearly dependent (Theorem 9).}

\textbf{Solution 13:} Column 3 is the zero vector. \textbf{Final Answer: Linearly dependent (Theorem 9).}

\textbf{Solution 14:} Check if one is a scalar multiple of the other. $-3\mathbf{v}_2 = -3 \begin{bmatrix} -1 \\ 2 \\ -3 \end{bmatrix} = \begin{bmatrix} 3 \\ -6 \\ 9 \end{bmatrix} = \mathbf{v}_1$. Since they are multiples, they are redundant. \textbf{Final Answer: Linearly dependent.}

\textbf{Solution 15:} There are only two vectors. $\mathbf{v}_1$ is not a scalar multiple of $\mathbf{v}_2$ (e.g., to match the first entry you need $c=2/5$, but $(2/5)(2) = 4/5 \neq 5$). \textbf{Final Answer: Linearly independent.}

\textbf{Solution 16:} False. The equation $A\mathbf{x} = \mathbf{0}$ \textit{always} has the trivial solution regardless of independence. For columns to be independent, it must have \textbf{only} the trivial solution.

\textbf{Solution 17:} False. Having $p \leq n$ just means Theorem 8 doesn't force dependence. They can still be dependent. For example, two vectors in $\mathbb{R}^3$ that are multiples of each other.

\textbf{Solution 18:} True. If the $5 \times 3$ matrix has 3 pivot columns, then every column is a pivot column. This means 0 free variables, which guarantees linear independence.

\textbf{Solution 19:} A $4 \times 2$ matrix with independent columns has exactly 2 columns, so it must have 2 pivots. Because each pivot steps down and to the right, they must be in rows 1 and 2. \textbf{Final Answer:}
$$ \begin{bmatrix} \blacksquare & * \\ 0 & \blacksquare \\ 0 & 0 \\ 0 & 0 \end{bmatrix} $$

\textbf{Solution 20:} A matrix has linearly independent columns if and only if there are no free variables in $A\mathbf{x} = \mathbf{0}$. Since there are 4 columns (4 variables), every single column must contain a pivot. \textbf{Final Answer: 4 pivot columns.}

\textbf{Solution 21:} $W(x) = \det\begin{bmatrix} e^{2x} & e^{3x} \\ 2e^{2x} & 3e^{3x} \end{bmatrix} = (e^{2x})(3e^{3x}) - (e^{3x})(2e^{2x}) = 3e^{5x} - 2e^{5x} = e^{5x}$. Since $e^{5x} \neq 0$, \textbf{Final Answer: Linearly independent.}

\textbf{Solution 22:} Setup the $3 \times 3$ Wronskian:
$$ W = \det\begin{bmatrix} 1 & x & e^x \\ 0 & 1 & e^x \\ 0 & 0 & e^x \end{bmatrix} $$
This is an upper triangular matrix. The determinant is the product of the diagonal: $1 \cdot 1 \cdot e^x = e^x$. Since $e^x \neq 0$, \textbf{Final Answer: Linearly independent.}

\textbf{Solution 23:} Recall the trigonometric identity $\cos^2(x) + \sin^2(x) = 1$. This means $\cos^2(x) + \sin^2(x) - 1(1) = 0$. We have found a non-trivial linear combination ($c_1=1, c_2=1, c_3=-1$) that equals the zero function. \textbf{Final Answer: Linearly dependent.}

\textbf{Solution 24:} The matrix is already upper triangular. $\det(A) = (1)(1)(5) = 5$. Since $\det(A) \neq 0$, the columns form a linearly independent set. \textbf{Final Answer: Linearly independent.}

\textbf{Solution 25:} Evaluate the determinant of the $3 \times 3$ matrix:
$$ \det(A) = k\begin{vmatrix} k & 1 \\ 1 & k \end{vmatrix} - 1\begin{vmatrix} 0 & 1 \\ 1 & k \end{vmatrix} + 0 = k(k^2 - 1) - 1(0 - 1) = k^3 - k + 1 $$
Wait, expanding correctly:
$\det(A) = k(k^2 - 1) - 1(0 - 1) + 0$? Let's re-evaluate:
Row 1: $k, 1, 0$.
$\det = k(k \cdot k - 1 \cdot 0) - 1(0 \cdot k - 1 \cdot 1) = k(k^2) - 1(-1) = k^3 + 1$.
For independence, $\det(A) \neq 0 \implies k^3 + 1 \neq 0 \implies k^3 \neq -1 \implies k \neq -1$.
\textbf{Final Answer: Independent for all real $k \neq -1$.}

\textbf{Solution 26:} We construct the payoff matrix $D = \begin{bmatrix} 10 & 10 & 0 \\ 20 & 10 & 10 \\ 30 & 10 & 20 \end{bmatrix}$. We solve $D\mathbf{x} = \mathbf{0}$:
$$ \begin{bmatrix} 10 & 10 & 0 \\ 20 & 10 & 10 \\ 30 & 10 & 20 \end{bmatrix} \xrightarrow{\substack{R_2 - 2R_1 \\ R_3 - 3R_1}} \begin{bmatrix} 10 & 10 & 0 \\ 0 & -10 & 10 \\ 0 & -20 & 20 \end{bmatrix} \xrightarrow{R_3 - 2R_2} \begin{bmatrix} 10 & 10 & 0 \\ 0 & -10 & 10 \\ 0 & 0 & 0 \end{bmatrix} $$
Column 3 lacks a pivot (free variable). The option can be perfectly replicated. \textbf{Final Answer: Linearly dependent (Arbitrage possible).}

\textbf{Solution 27:} Yes. The 3 assets' payoffs are modeled as vectors in $\mathbb{R}^2$ (since there are 2 states). We have $p = 3$ vectors in $n = 2$ dimensions. By Theorem 8, any set where $p > n$ is linearly dependent. \textbf{Final Answer: Yes, by Theorem 8, at least one asset is redundant.}

\textbf{Solution 28: The Flaw:} The student confused the number of vectors ($p$) with the number of entries/dimensions ($n$). For 3 vectors in $\mathbb{R}^5$, $p=3$ and $n=5$. Since $p < n$, Theorem 8 does not apply. \textbf{The Correction:} The set is not automatically dependent; one must perform row reduction to determine independence.

\textbf{Solution 29: The Flaw:} The student analyzed the rows instead of the columns. Linear independence requires a pivot in every \textit{column}. \textbf{The Correction:} The matrix has 3 columns but only 2 pivots. Therefore, column 3 represents a free variable. The system has non-trivial solutions, making the vectors linearly dependent.

\textbf{Solution 30: The Flaw:} Setting all weights to zero yields the trivial solution, which always works for any set of vectors. \textbf{The Correction:} Linear independence requires the trivial solution to be the \textit{only} solution. Merely verifying that the trivial solution works proves nothing; you must show no other solutions exist.

\newpage
\section{Concept Checklist \& Question Mapping}
This checklist serves as a comprehensive inventory of all linear independence concepts, tricks, and problem types extracted from the source material. It maps each concept to the relevant questions in the problem set to ensure total coverage.

\begin{itemize}
    \item \textbf{C1: Algorithmic Reduction (Row Reduction / Gaussian Elimination)} 
    \newline \textit{Questions:} 1, 2, 3
    
    \item \textbf{C2: Independence of Matrix Columns}
    \newline \textit{Questions:} 4, 5
    
    \item \textbf{C3: Unknown Variables and Spans}
    \newline \textit{Questions:} 6, 7, 8, 9
    
    \item \textbf{C4: Visual Inspection - Oversize Sets (Theorem 8)}
    \newline \textit{Questions:} 10, 11
    
    \item \textbf{C5: Visual Inspection - Zero Vector Rule (Theorem 9)}
    \newline \textit{Questions:} 12, 13
    
    \item \textbf{C6: Visual Inspection - Two-Vector Rule}
    \newline \textit{Questions:} 14, 15
    
    \item \textbf{C7: Conceptual True/False}
    \newline \textit{Questions:} 16, 17, 18
    
    \item \textbf{C8: Echelon Form Analysis}
    \newline \textit{Questions:} 19, 20
    
    \item \textbf{C9: Advanced Applications - Linear Independence of Functions (Wronskian)}
    \newline \textit{Questions:} 21, 22, 23
    
    \item \textbf{C10: Advanced Applications - Independence via Determinants}
    \newline \textit{Questions:} 24, 25
    
    \item \textbf{C11: Real-World Modeling \& Quantitative Finance}
    \newline \textit{Questions:} 26, 27
    
    \item \textbf{C12: Diagnostic Testing ("Find the Flaw")}
    \newline \textit{Questions:} 28, 29, 30
\end{itemize}

\end{document}